Este capítulo cubre el análisis del sistema de información a desarrollar, haciendo uso del lenguaje de modelado UML.

\section{Modelo Conceptual}
A partir de los requisitos de información, se desarrollará un diagrama conceptual de clases UML, identificando las clases, atributos, relaciones, restricciones adicionales y reglas de derivación necesarias.

\section{Modelo de Casos de Uso}
A partir de los requisitos funcionales descritos anteriormente, se emplearan los casos de uso como mecanismo para representar las interacciones entre los actores y el sistema bajo estudio. Para cada caso de uso deberá indicarse los actores implicados, las precondiciones y postcondiciones, los pasos que conforman el escenario principal y el conjunto de posibles escenarios alternativos. Los actores pueden ser roles de personas físicas, sistemas externos o incluso el tiempo (para eventos temporales).

\section{Modelo de Interfaz de Usuario}
En esta sección se deberá incluir un prototipo de baja fidelidad o mockup de la interfaz de usuario del sistema. Además, es preciso elaborar un diagrama de navegación, reflejando la secuencia de pantallas a las que tienen acceso los diferentes roles de usuario y la conexión entre éstas.

\section{Modelo de Comportamiento}
En esta sección opcional, se incluirá el modelo de comportamiento desarrollado a partir de los casos de uso definidos. Para ello, se realizarán los diagramas de secuencia del sistema, donde se identificarán las operaciones o servicios del sistema. Luego, se detallará el contrato de las operaciones identificadas.
